\section{推荐系统}
\label{sec:chapter14}

\subsection{基于内容的推荐算法}

以向用户推荐电影为例。假设已经为每部电影构建了特征向量，例如不同类型或题材所占的程度，并且已知每个用户对部分电影的评分。对于某个用户，可以将电影特征作为自变量 $x$、该用户的评分作为因变量 $y$，从而为每个用户建立一个独立的\textbf{线性回归}模型。

设共有 $n_u$ 个用户和 $n_m$ 部电影，第 $i$ 部电影的特征向量为 $x^{(i)}\in\mathbb R^{n+1}$（包含偏置项）。$r(i,j)$ 表示用户 $j$ 是否评价了电影 $i$；若已经评分，则对应分数记为 $y^{(i,j)}$. 对于用户 $j$，只使用他已经评价过的电影训练参数 $\theta^{(j)}$，优化目标为：
\[
\min_{\theta^{(j)}}~~\frac{1}{2}\sum_{i:r(i,j)=1}\left((\theta^{(j)})^Tx^{(i)}-y^{(i,j)}\right)^2+\frac{\lambda}{2}\sum_{k=1}^{n}(\theta_k^{(j)})^2
\]
每个用户的线性回归在形式上相互独立，但可以把所有用户的目标函数相加，统一训练全部参数：
\[
J(\theta^{(1)},\ldots,\theta^{(n_u)}):=\frac{1}{2}\sum_{j=1}^{n_u}\sum_{i:r(i,j)=1}\left((\theta^{(j)})^Tx^{(i)}-y^{(i,j)}\right)^2+\frac{\lambda}{2}\sum_{j=1}^{n_u}\sum_{k=1}^{n}(\theta_k^{(j)})^2
\]
为了使用梯度下降等方法优化，所需的偏导数为：
\[
\frac{\partial J}{\partial\theta^{(j)}_k}=\sum_{i:r(i,j)=1}\left((\theta^{(j)})^Tx^{(i)}-y^{(i,j)}\right)\theta^{(j)}_k+\lambda\theta_k^{(j)}[k>0]
\]
训练完成后，对于特征向量为 $x$ 的任意电影，用户 $j$ 的预测评分为 $(\theta^{(j)})^Tx$. 可以据此对用户尚未评价的电影排序，并推荐预测评分较高的项目。

基于内容的推荐算法有一个明显限制：必须预先知道每部电影的特征向量。然而，哪些特征最能解释用户偏好往往难以确定，人工标注大量电影也需要较高成本。为避免依赖显式内容特征，可以使用协同过滤，从用户评分数据中同时学习电影和用户的表示。


\subsection{协同过滤算法}

\paragraph{初始版本}

现在假设电影特征未知，但可以通过询问用户等方式获得用户参数 $\theta^{(j)}$，例如用户对不同类型电影的偏好。这样便可以反过来把用户参数作为已知量，根据评分学习电影特征 $x^{(i)}$. 对于第 $i$ 部电影，只使用评价过该电影的用户，优化目标为：
\[
\min_{\theta^{(j)}}~~\frac{1}{2}\sum_{j:r(i,j)=1}\left((\theta^{(j)})^Tx^{(i)}-y^{(i,j)}\right)^2+\frac{\lambda}{2}\sum_{k=1}^{n}(x_k^{(i)})^2
\]
与前面按用户训练参数类似，各电影特征的优化也相互独立，可以将它们的目标相加后统一写为：
\[
J(x^{(1)},\ldots,x^{(n_m)}):=\frac{1}{2}\sum_{i=1}^{n_m}\sum_{j:r(i,j)=1}\left((\theta^{(j)})^Tx^{(i)}-y^{(i,j)}\right)^2+\frac{\lambda}{2}\sum_{i=1}^{n_m}\sum_{k=1}^{n}(x_k^{(i)})^2
\]
为了优化电影特征，所需的偏导数为：
\[
\frac{\partial J}{\partial x^{(i)}_k}=\sum_{j:r(i,j)=1}\left((\theta^{(j)})^Tx^{(i)}-y^{(i,j)}\right)\theta^{(i)}_k+\lambda x_k^{(i)}[k>0]
\]

总结而言，已知 $\theta^{(j)}$ 时可以学习 $x^{(i)}$，已知 $x^{(i)}$ 时也可以学习 $\theta^{(j)}$. 因此，可以先随机初始化一组用户参数，用它们学习所有电影特征；再固定新得到的电影特征，重新学习用户参数；随后继续使用新的用户参数更新电影特征。如此在两组参数之间\textbf{交替优化}，直至电影特征和用户参数都趋于稳定，这构成了协同过滤算法的初始版本。


\paragraph{改进版本}

实际上，没有必要把两组参数拆开并手工反复交替训练。观察“固定用户参数学习电影特征”和“固定电影特征学习用户参数”这两个优化目标，可以发现它们的非正则化项完全相同，均为 $\sum_{(i,j):r(i,j)=1}\left((\theta^{(j)})^Tx^{(i)}-y^{(i,j)}\right)^2$. 因此，只需合并两边的正则化项，就可以把电影特征和用户参数放入同一个目标函数中同时优化：
\begin{gather*}
J(x^{(1)},\ldots,x^{(n_m)},\theta^{(1)},\ldots,\theta^{(n_u)})=\\\frac{1}{2}\sum_{(i,j):r(i,j)=1}\left((\theta^{(j)})^Tx^{(i)}-y^{(i)}\right)^2+\frac{\lambda}{2}\sum_{i=1}^{n_m}\sum_{k=1}^n(x^{(i)}_k)^2+\frac{\lambda}{2}\sum_{j=1}^{n_u}\sum_{k=1}^n(\theta^{(j)}_k)^2
\end{gather*}
值得注意的是，在改进版本中，特征由模型自动学习。由于电影特征和用户参数可以共同吸收整体偏移，模型中\textbf{无须额外加入偏置项}，因此 $x^{(i)}\in\mathbb R^n,\theta^{(j)}\in\mathbb R^n$.

上式的偏导数为：
\begin{align*}
&\frac{\partial J}{\partial\theta^{(j)}_k}=\sum_{i:r(i,j)=1}\left((\theta^{(j)})^Tx^{(i)}-y^{(i,j)}\right)\theta^{(j)}_k+\lambda\theta_k^{(j)}\\
&\frac{\partial J}{\partial x^{(i)}_k}=\sum_{j:r(i,j)=1}\left((\theta^{(j)})^Tx^{(i)}-y^{(i,j)}\right)x^{(i)}_k+\lambda x_k^{(i)}
\end{align*}
得到目标函数和梯度后，可以使用梯度下降或其他算法（如 LBFGS 等）进行优化。


\paragraph{向量化版本}

逐项遍历所有电影和用户会产生较大开销。为提高代码运行效率，需要将协同过滤算法写成矩阵形式并进行并行计算。

构造评分矩阵 $Y:=\begin{bmatrix}y^{(i,j)}\end{bmatrix}\in\mathbb R^{n_m\times n_u}$，其第 $i$ 行第 $j$ 列表示用户 $j$ 对电影 $i$ 的评分；构造电影特征矩阵 $X:=\begin{bmatrix}(x^{(1)})^T\\ \vdots\\ (x^{(n_m)})^T\end{bmatrix}\in\mathbb R^{n_m\times n}$，其第 $i$ 行是电影 $i$ 的特征向量；再构造用户参数矩阵 $\Theta:=\begin{bmatrix}(\theta^{(1)})^T\\ \vdots\\ (\theta^{(n_u)})^T\end{bmatrix}\in\mathbb R^{n_u\times n}$，其第 $j$ 行是用户 $j$ 的参数向量。由此，所有用户对所有电影的预测评分可以一次写成矩阵：
\[
X\Theta^T\in\mathbb R^{n_m\times n_u}
\]
利用 \texttt{numpy} 的语法可以简洁地实现代价函数及其梯度的向量化计算，详见代码。


\subsection{代码实现}

\begin{lstlisting}[language=Python]
import numpy as np
from scipy.io import loadmat
from scipy.optimize import minimize

def J(Y, R, X, Theta, lamb):
	return 0.5 * (np.sum(((X @ Theta.T - Y) * R) ** 2) + \
		lamb * np.sum(X ** 2) + lamb * np.sum(Theta ** 2))

def partJ(Y, R, X, Theta, lamb):
	return np.concatenate((
		( ((X @ Theta.T - Y) * R) @ Theta + lamb * X ).flatten(), 
		( ((X @ Theta.T - Y) * R).T @ X + lamb * Theta ).flatten()
	))

def train(Y, R, lamb, n):
	(n_m, n_u) = Y.shape
	xt = np.empty(n*n_m+n*n_u)
	return minimize(fun = lambda xt, Y, R, lamb : J(Y, R, xt[:n*n_m].reshape((n_m, n)), \
													xt[n*n_m:].reshape((n_u, n)), lamb), 
					x0 = np.random.randn(n*n_m+n*n_u), 
					args = (Y, R, lamb), 
					method = 'TNC', 
					jac = lambda xt, Y, R, lamb: partJ(Y, R, xt[:n*n_m].reshape((n_m, n)), \
													xt[n*n_m:].reshape((n_u, n)), lamb)
					)

def predict(Y, R, xt, n):
	(n_m, n_u) = Y.shape
	X, Theta = xt[:n*n_m].reshape((n_m, n)), xt[n*n_m:].reshape((n_u, n))
	return X @ Theta.T

data = loadmat('ex8_movies.mat')
Y = data['Y']
R = data['R']
Ymean = Y.mean(axis=1, keepdims=True)
res = train(Y-Ymean, R, lamb=1, n=50)
print(res)
np.save('xt.npy', res.x)
xt = res.x
\end{lstlisting}

优化结果为：

\begin{lstlisting}
     fun: 11078.825074101991
     jac: array([ 4.75273039e-07, -8.33595791e-07, -4.85091646e-07, ...,
        9.67607422e-07,  3.80539749e-06,  1.86386969e-06])
 message: 'Converged (|f_n-f_(n-1)| ~= 0)'
    nfev: 14671
     nit: 463
  status: 1
 success: True
       x: array([ 0.30269431, -1.3577984 , -0.19821756, ...,  0.12718836,
       -0.40793964, -0.60753772])
\end{lstlisting}

优化得到电影特征和用户参数后，先计算完整的预测评分矩阵，再从中选取第一个用户预测评分最高的 $10$ 部电影：

\begin{lstlisting}[language=Python]
xt = np.load('xt.npy')
pred = predict(Y, R, xt, n=50) + Ymean
movie_list = []
with open('movie_ids.txt', encoding='latin-1') as file:
	for line in file:
		movie_list.append(' '.join(line.strip().split(' ')[1: ]))
movie_list = np.array(movie_list)
idx = np.argsort(pred[:, 0])[::-1]
print('Top 10 movies for user 1:')
for movie in movie_list[idx][:10]:
	print(movie)
\end{lstlisting}

结果为：

\begin{lstlisting}
Top 10 movies for user 1:
Titanic (1997)
In the Name of the Father (1993)
Philadelphia (1993)
Duck Soup (1933)
Ice Storm, The (1997)
Saint, The (1997)
William Shakespeare's Romeo and Juliet (1996)
Boot, Das (1981)
People vs. Larry Flynt, The (1996)
Manhattan (1979)
\end{lstlisting}
